% =====================================================================
%  Guía para construir una propuesta de ALTO VALOR
%  Samsung Solve for Tomorrow 2026
%  I.E. Sor María Juliana · Cartago, Valle del Cauca
%  PhD. Álvaro Cárdenas Orozco · Plataforma ConectaTE
%  Motor: XeLaTeX
% =====================================================================
\documentclass[11pt, letterpaper]{article}

\usepackage[spanish, es-noindentfirst]{babel}
\usepackage{geometry}
\geometry{top=2.3cm, bottom=2.2cm, left=2.3cm, right=2.3cm, headheight=15pt}

\usepackage{xcolor}
\definecolor{smjAzul}{HTML}{0D1B3E}
\definecolor{smjVino}{HTML}{6B1A2A}
\definecolor{smjDorado}{HTML}{C9A84C}
\definecolor{smjGris}{HTML}{F4F4F4}
\definecolor{sftAzul}{HTML}{1428A0}     % azul Samsung (acento del programa)
\definecolor{valorVerde}{HTML}{1A7A4C}  % verde "alto valor"

\usepackage{fontspec}
\setmainfont{texgyretermes}[Extension=.otf, UprightFont=*-regular, BoldFont=*-bold, ItalicFont=*-italic, BoldItalicFont=*-bolditalic]
\setsansfont{texgyreheros}[Extension=.otf, UprightFont=*-regular, BoldFont=*-bold, ItalicFont=*-italic, BoldItalicFont=*-bolditalic]
\newcommand{\sft}{\textsf}

\usepackage{microtype}
\usepackage{amssymb}
\usepackage{graphicx}
\usepackage{tikz}
\usetikzlibrary{positioning, arrows.meta, calc}
\usepackage{pgffor}
\usepackage{enumitem}
\setlist{leftmargin=1.4em, itemsep=2pt, topsep=3pt}
\usepackage{booktabs, tabularx, array, colortbl}
\usepackage{multicol}
\usepackage{fancyhdr}
\usepackage{titlesec}
\usepackage{parskip}
\usepackage[hidelinks]{hyperref}
\hypersetup{colorlinks=true, linkcolor=smjAzul, urlcolor=sftAzul,
  pdfauthor={PhD. Álvaro Cárdenas Orozco},
  pdftitle={Guía de propuestas de alto valor — Solve for Tomorrow 2026}}

% --- Encabezado / pie institucional ---
\pagestyle{fancy}
\fancyhf{}
\fancyhead[L]{\footnotesize\sft{\textcolor{smjAzul}{Solve for Tomorrow 2026}}}
\fancyhead[R]{\footnotesize\sft{\textcolor{smjVino}{I.E. Sor María Juliana}}}
\fancyfoot[C]{\footnotesize\sft{\textcolor{gray}{Guía de propuestas de alto valor \;·\; ConectaTE \;·\; PhD. Álvaro Cárdenas Orozco \;·\; \thepage}}}
\renewcommand{\headrulewidth}{0.4pt}
\renewcommand{\headrule}{\hbox to\headwidth{\color{smjDorado}\leaders\hrule height \headrulewidth\hfill}}
\renewcommand{\footrulewidth}{0.2pt}

% --- Estilo de secciones ---
\titleformat{\section}{\sffamily\large\bfseries\color{smjAzul}}{\thesection.}{0.5em}{}[{\color{smjDorado}\titlerule[1pt]}]
\titleformat{\subsection}{\sffamily\normalsize\bfseries\color{smjVino}}{\thesubsection}{0.5em}{}
\titlespacing{\section}{0pt}{14pt}{6pt}

% --- Cajas educativas ---
\usepackage{tcolorbox}
\tcbuselibrary{skins, breakable}

\newtcolorbox{oficial}{enhanced, breakable, colback=sftAzul!5, colframe=sftAzul, boxrule=1pt, arc=2pt,
  coltitle=white, colbacktitle=sftAzul, fonttitle=\sffamily\bfseries\small, title={Dato oficial verificado}}
\newtcolorbox{clavebox}{enhanced, breakable, colback=valorVerde!7, colframe=valorVerde, boxrule=0.9pt, arc=2pt,
  coltitle=white, colbacktitle=valorVerde, fonttitle=\sffamily\bfseries\small, title={Claves de alto valor}}
\newtcolorbox{evitabox}{enhanced, breakable, colback=smjVino!6, colframe=smjVino, boxrule=0.9pt, arc=2pt,
  coltitle=white, colbacktitle=smjVino, fonttitle=\sffamily\bfseries\small, title={Errores que restan valor}}
\newtcolorbox{ejemplobox}{enhanced, breakable, colback=smjGris, colframe=smjDorado!85!black, boxrule=0.9pt, arc=2pt,
  coltitle=smjAzul, colbacktitle=smjDorado!45, fonttitle=\sffamily\bfseries\small, title={Ejemplo modelo · contexto Cartago (proyecto «AquaVigía»)}}
\newtcolorbox{milcbox}{enhanced, breakable, colback=smjDorado!9, colframe=smjDorado!85!black, boxrule=0.9pt, arc=2pt,
  coltitle=smjAzul, colbacktitle=smjDorado!40, fonttitle=\sffamily\bfseries\small, title={La lente MILC}}
\newtcolorbox{pidebox}[1]{enhanced, breakable, colback=smjAzul!4, colframe=smjAzul, boxrule=1.1pt, arc=2pt,
  coltitle=white, colbacktitle=smjAzul, fonttitle=\sffamily\bfseries, title={#1}}
\newtcolorbox{borradorbox}[1]{enhanced, breakable, colback=white, colframe=smjAzul!45, boxrule=0.6pt, arc=2pt,
  coltitle=smjAzul, colbacktitle=smjGris, fonttitle=\sffamily\bfseries\small, title={Tu borrador — #1}}

\newcommand{\renglones}[1]{\foreach \n in {1,...,#1}{\par\vspace{0.52cm}\textcolor{smjAzul!22}{\rule{\linewidth}{0.3pt}}}}
\newcommand{\chk}{\textcolor{smjAzul}{$\square$}~}
\newcommand{\preg}[1]{\textit{\textcolor{smjVino}{#1}}}

\begin{document}

% =====================================================================
% PORTADA
% =====================================================================
\begin{titlepage}
\thispagestyle{empty}
\begin{center}
  {\color{smjAzul}\rule{\linewidth}{2pt}}\\[0.25cm]
  {\large\sffamily\textsc{\color{smjAzul} Institución Educativa Sor María Juliana}}\\[0.1cm]
  {\normalsize\sffamily\color{smjVino} Cartago, Valle del Cauca \;·\; Área de Tecnología e Informática}\\[0.1cm]
  {\small\sffamily\color{smjDorado!80!black} Plataforma educativa ConectaTE}\\[0.25cm]
  {\color{smjDorado}\rule{\linewidth}{1pt}}\\[1.6cm]

  {\small\sffamily\color{sftAzul}\bfseries GUÍA DE PREPARACIÓN \;·\; INNOVACIÓN SOCIAL STEM}\\[0.5cm]
  {\Huge\sffamily\bfseries\color{smjAzul} Construye una propuesta\\[0.2cm] de \textcolor{valorVerde}{alto valor}}\\[0.7cm]
  {\large\sffamily\color{smjVino} Samsung \textit{Solve for Tomorrow} 2026}\\[0.3cm]
  {\normalsize\color{smjAzul!80} Del problema de tu comunidad a una solución STEM ganadora.}\\[1.4cm]

  \begin{tikzpicture}[font=\sffamily\footnotesize\bfseries]
    \foreach \i/\num/\n/\c in {0/1/Empatizar/smjAzul, 1/2/Definir/smjVino, 2/3/Idear/smjDorado!80!black, 3/4/Prototipar/valorVerde, 4/5/Validar/sftAzul} {
      \node[draw=\c, fill=\c!12, rounded corners=3pt, minimum width=2.05cm, minimum height=1.05cm, align=center, text=\c] at (\i*2.45,0) {\num\\\n};
      \ifnum\i>0 \draw[-{Stealth[length=2mm]}, smjDorado, line width=1pt] (\i*2.45-1.25,0) -- (\i*2.45-0.85,0);\fi
    }
  \end{tikzpicture}\\[0.2cm]
  {\footnotesize\sffamily\color{gray} El método de diseño que evalúa el formulario oficial}\\[1.8cm]

  {\color{smjDorado}\rule{0.5\linewidth}{0.5pt}}\\[0.6cm]
  \begin{minipage}{0.46\textwidth}\begin{flushleft}
    \textbf{\sffamily\color{smjAzul} Equipo:}\\ \textcolor{gray}{Nombre del equipo y grados}
  \end{flushleft}\end{minipage}\hfill
  \begin{minipage}{0.46\textwidth}\begin{flushright}
    \textbf{\sffamily\color{smjAzul} Docente guía:}\\ PhD. Álvaro Cárdenas Orozco
  \end{flushright}\end{minipage}\\[1.6cm]

  {\footnotesize\color{gray} Material de preparación independiente. \emph{Solve for Tomorrow} es un programa de Samsung Electronics; la inscripción es exclusivamente oficial.}\\[0.3cm]
  {\color{smjAzul}\rule{\linewidth}{2pt}}
\end{center}
\end{titlepage}

% =====================================================================
\section{Antes de empezar: ¿qué es y quién puede participar?}

\emph{Solve for Tomorrow} es el programa de ciudadanía corporativa de Samsung que reta a estudiantes de \textbf{colegios públicos} a diseñar una \textbf{solución tecnológica a un problema real de su comunidad}. No premia la idea más vistosa, sino la que \textbf{entiende de verdad un problema} y propone una salida \textbf{aplicable y con sentido}. Por eso esta guía no te enseña a llenar casillas: te enseña a pensar como un equipo que transforma su territorio.

\begin{oficial}
\setlength{\parskip}{3pt}
\textbf{Requisitos (edición 2026, la décima en Colombia).}
\begin{itemize}[nosep]
  \item Estudiantes de \textbf{colegios públicos} que cursen \textbf{8.\textsuperscript{o} a 12.\textsuperscript{o}}, con edades entre \textbf{13 y 19 años}.
  \item Equipo de \textbf{2 a 4 estudiantes} (revisa en las bases la exigencia de equipos mixtos) acompañados por \textbf{un docente} del mismo colegio.
  \item Enfoque \textbf{STEM} (ciencia, tecnología, ingeniería, matemática) trabajado con la metodología \textbf{Design Thinking}.
\end{itemize}
\textbf{Trayectoria del programa:} +48.000 estudiantes, 2.353 colegios públicos y cerca de 14.200 proyectos en los 32 departamentos del país.

\smallskip
\textbf{Importante:} este documento es una \textbf{guía de trabajo}. \textbf{No es un medio de inscripción.} La inscripción se hace únicamente en el sitio oficial, en las fechas vigentes: \url{https://www.samsung.com/co/campaign/solvefortomorrow/}
\end{oficial}

% =====================================================================
\section{¿Qué hace que una propuesta sea de alto valor?}

Los jurados leen cientos de proyectos. Los que avanzan comparten cinco rasgos. Tenlos como brújula en \emph{cada} campo que escribas.

\begin{center}
\renewcommand{\arraystretch}{1.35}
\begin{tabularx}{\linewidth}{@{}p{0.32\linewidth} X@{}}
\toprule
\textbf{\color{smjAzul}Rasgo} & \textbf{\color{smjAzul}Qué significa en la práctica}\\
\midrule
\textbf{1. Problema real y específico} & No «la contaminación» en abstracto, sino «la quebrada del barrio X se desborda con basura cada invierno». Concreto, ubicado, verificable.\\
\textbf{2. Empatía con evidencia} & Hablaste con la comunidad. Tienes datos, frases, observaciones. No supones lo que la gente necesita: lo \emph{sabes}.\\
\textbf{3. Solución diferenciada} & Explicas por qué tu idea es mejor o distinta de lo que ya existe. Si ya hay una app igual, ¿qué aportas tú?\\
\textbf{4. STEM aplicado y explícito} & Se ve la ciencia, la tecnología, la ingeniería o la matemática \emph{funcionando} dentro de la solución, no como adorno.\\
\textbf{5. Impacto medible} & Dices cómo sabrás si funcionó: una cifra, un antes y un después, una forma de validarlo con la comunidad.\\
\bottomrule
\end{tabularx}
\end{center}

\begin{milcbox}
Este concurso es el modelo \textbf{MILC} en acción: \emph{investigar el territorio y crear tecnología con sentido}. Un buen proyecto Solve for Tomorrow nace de mirar tu comunidad con las cinco dimensiones del MILC (personal, emocional, ciudadana, local e intergeneracional). Volveremos a esto al final para darle profundidad a tu propuesta.
\end{milcbox}

% =====================================================================
\section{El mapa de tu propuesta: Design Thinking en 5 etapas}

La sección \textbf{«Propuesta»} del formulario es el corazón evaluable, y sigue las cinco etapas del \textbf{Design Thinking}. Cada una tiene un \textbf{límite de caracteres}: escribir bien con poco espacio es parte del reto. Trabaja primero en sucio (aquí) y pule hasta que cada palabra cuente.

A lo largo de las cinco etapas seguiremos un mismo ejemplo, \textbf{«AquaVigía»}, para que veas un hilo coherente de principio a fin. Es un caso \emph{ilustrativo}: tu proyecto será el tuyo.

% ---------- ETAPA 1 ----------
\subsection{Etapa 1 · Empatizar}
\begin{pidebox}{¿Cuál es su nivel de empatía con la comunidad? \hfill\normalfont\small máx. 840 caracteres}
Identifiquen con claridad la \textbf{comunidad objetivo}, pónganse en su lugar y describan de forma concreta sus \textbf{necesidades, deseos y desafíos}.
\smallskip\\
\preg{Preguntas de apoyo: ¿Quién o cuál es la comunidad objetivo? ¿Cuáles son sus características? ¿Dónde está ubicada? ¿Por qué es importante para ustedes? ¿Cómo se siente la comunidad sobre su situación actual? ¿Cuáles son sus expectativas de cambio?}
\end{pidebox}

\begin{clavebox}
\begin{itemize}[nosep]
  \item \textbf{Nombra y ubica} a la comunidad: barrio, vereda, edad, número aproximado de personas.
  \item Incluye \textbf{una voz real}: una frase que alguien te dijo, una escena que observaste. Eso demuestra empatía verdadera.
  \item Describe lo que \textbf{sienten}, no solo lo que les falta. La emoción mueve al jurado y revela que estuviste ahí.
\end{itemize}
\end{clavebox}
\begin{evitabox}
\begin{itemize}[nosep]
  \item Hablar de «la sociedad» o «los jóvenes» en general: demasiado amplio, no se siente real.
  \item Suponer necesidades sin haber preguntado nunca a nadie de la comunidad.
\end{itemize}
\end{evitabox}
\begin{ejemplobox}
Las 40 familias del sector La Platanera, junto a la quebrada Las Cruces en Cartago, usan el agua de la quebrada para lavar y, a veces, para cocinar. En invierno baja turbia y con basura; varios niños han tenido diarrea. Doña Ruby nos dijo: «uno no sabe si está limpia hasta que ya está enfermo». La comunidad quiere una forma sencilla y barata de saber, día a día, si el agua es segura.
\end{ejemplobox}
\begin{borradorbox}{máx. 840 caracteres (≈ 6 renglones)}
\renglones{6}
\end{borradorbox}

% ---------- ETAPA 2 ----------
\subsection{Etapa 2 · Definir}
\begin{pidebox}{¿Cuál es el problema identificado? \hfill\normalfont\small máx. 840 caracteres}
Planteen con claridad el \textbf{problema que buscan resolver}, basándose en lo que comprendieron de la comunidad.
\smallskip\\
\preg{Preguntas de apoyo: ¿Cuál es el problema que afecta a la comunidad? ¿Qué impacto tiene en su vida y en el entorno? ¿Por qué es importante abordarlo?}
\end{pidebox}
\begin{clavebox}
\begin{itemize}[nosep]
  \item Redacta el problema como \textbf{una frase clara} (causa $\rightarrow$ efecto): «Como X, entonces Y».
  \item Muestra el \textbf{impacto}: a cuántos afecta, qué consecuencias tiene (salud, dinero, tiempo, ambiente).
  \item El problema debe \textbf{nacer de la etapa 1}. Si no se conecta con lo que empatizaste, pierde fuerza.
\end{itemize}
\end{clavebox}
\begin{evitabox}
\begin{itemize}[nosep]
  \item Confundir el problema con la falta de tu solución («el problema es que no tienen una app»). El problema existe antes que tu idea.
  \item Plantear un problema tan grande que ningún equipo escolar podría abordarlo.
\end{itemize}
\end{evitabox}
\begin{ejemplobox}
Las familias de La Platanera no tienen forma de saber si el agua de la quebrada es segura antes de usarla. Como el deterioro no siempre se ve a simple vista, la usan «a la suerte», y eso deriva en enfermedades estomacales, sobre todo en los niños, y en gastos médicos que estas familias de bajos ingresos no pueden asumir. Es urgente porque el agua se usa todos los días.
\end{ejemplobox}
\begin{borradorbox}{máx. 840 caracteres (≈ 6 renglones)}
\renglones{6}
\end{borradorbox}

% ---------- ETAPA 3 ----------
\subsection{Etapa 3 · Idear}
\begin{pidebox}{¿Cuál es la idea de solución que proponen? \hfill\normalfont\small máx. 560 caracteres}
Presenten con claridad la \textbf{solución} y justifiquen por qué resuelve el problema.
\smallskip\\
\preg{Preguntas de apoyo: ¿Cuál es la solución? ¿Por qué resuelve el problema identificado? ¿Cómo se diferencia de las soluciones que ya existen en la comunidad?}
\end{pidebox}
\begin{clavebox}
\begin{itemize}[nosep]
  \item Di la solución \textbf{en una frase} antes de explicarla. El jurado debe entenderla en 5 segundos.
  \item Responde el «\textbf{por qué resuelve}»: une la solución con la causa del problema de la etapa 2.
  \item El campo \textbf{exige} decir en qué te diferencias. Ten lista esa respuesta: precio, sencillez, que la maneje la propia comunidad, etc.
\end{itemize}
\end{clavebox}
\begin{evitabox}
\begin{itemize}[nosep]
  \item Proponer «una app» sin más: es lo que todos dicen. Sé específico sobre qué hace y para quién.
  \item Olvidar la diferenciación. Si no dices por qué tu idea es distinta, parece repetida.
\end{itemize}
\end{evitabox}
\begin{ejemplobox}
\textbf{AquaVigía}: un medidor comunitario de bajo costo que revisa la turbidez y el pH del agua y enciende una luz verde, amarilla o roja, fácil de leer por cualquiera. Resuelve el problema porque convierte algo invisible (¿está segura?) en una señal clara y diaria. A diferencia de un laboratorio (caro y lejano), lo opera la propia comunidad y cuesta poco.
\end{ejemplobox}
\begin{borradorbox}{máx. 560 caracteres (≈ 4 renglones)}
\renglones{4}
\end{borradorbox}

% ---------- ETAPA 4 ----------
\subsection{Etapa 4 · Prototipar (y dónde vive el STEM)}
\begin{pidebox}{¿Cómo funciona la solución propuesta? \hfill\normalfont\small máx. 840 caracteres}
Expliquen cómo funcionaría la solución y \textbf{cómo aplican STEM} en ella, describiendo un modelo inicial de prototipo. \emph{No necesitan tener el prototipo físico: un prototipo descriptivo (en palabras) es válido.}
\smallskip\\
\preg{Preguntas de apoyo: ¿Cómo funciona su solución? (describan el proceso) ¿Qué elementos tiene el prototipo? (materiales, funcionalidades) ¿Cómo se relaciona con la ciencia, la tecnología, la ingeniería y/o las matemáticas?}
\end{pidebox}
\begin{clavebox}
\begin{itemize}[nosep]
  \item Describe el \textbf{funcionamiento paso a paso}: qué entra, qué hace, qué sale.
  \item \textbf{Nombra el STEM con todas las letras.} Es lo que más distingue a este concurso: di qué ciencia, qué tecnología, qué ingeniería y qué matemática usas (ver la Sección 5).
  \item Lista \textbf{elementos concretos}: sensores, materiales, app, etc. Lo concreto se siente factible.
\end{itemize}
\end{clavebox}
\begin{evitabox}
\begin{itemize}[nosep]
  \item Describir la solución como magia («detecta todo automáticamente») sin explicar el cómo.
  \item Mencionar STEM como etiqueta sin mostrar dónde actúa de verdad.
\end{itemize}
\end{evitabox}
\begin{ejemplobox}
Un sensor de turbidez y otro de pH (conectados a una placa programable tipo micro:bit/Arduino) toman lecturas; el código las compara con umbrales seguros y enciende un LED verde/amarillo/rojo. \textbf{Ciencia}: rangos de pH y turbidez aptos para contacto. \textbf{Tecnología}: sensores + microcontrolador. \textbf{Ingeniería}: una caja resistente al agua y de fácil mantenimiento. \textbf{Matemática}: promedios diarios y umbrales para decidir el color.
\end{ejemplobox}
\begin{borradorbox}{máx. 840 caracteres (≈ 6 renglones)}
\renglones{6}
\end{borradorbox}

% ---------- ETAPA 5 ----------
\subsection{Etapa 5 · Validar}
\begin{pidebox}{¿Cuál es su visión sobre el impacto positivo en la comunidad? \hfill\normalfont\small máx. 560 caracteres}
Argumenten cómo \textbf{probarán} la solución en la comunidad y proyecten sus \textbf{efectos positivos}.
\smallskip\\
\preg{Preguntas de apoyo: ¿Cómo contactarán a la comunidad para recibir retroalimentación? ¿Cuál sería el impacto esperado? ¿Cómo medirían el éxito? ¿Cómo involucrarían a la comunidad en la implementación?}
\end{pidebox}
\begin{clavebox}
\begin{itemize}[nosep]
  \item Da \textbf{una forma de medir el éxito}: una cifra o un antes/después («reducir las diarreas», «que 30 familias lo usen»).
  \item Muestra que la comunidad \textbf{participa}, no solo recibe: se apropian, lo mantienen, lo cuidan.
  \item Sé realista: un plan creíble vale más que una promesa enorme.
\end{itemize}
\end{clavebox}
\begin{evitabox}
\begin{itemize}[nosep]
  \item Decir «tendrá un gran impacto» sin explicar cómo lo sabrías.
  \item Dejar a la comunidad como espectadora pasiva de tu solución.
\end{itemize}
\end{evitabox}
\begin{ejemplobox}
Instalaríamos un prototipo con la junta de acción comunal y formaríamos a dos familias para leerlo y reportar. Mediríamos éxito con dos datos: número de familias que consultan el medidor antes de usar el agua, y casos de enfermedad estomacal reportados en el puesto de salud, comparando tres meses antes y después. La comunidad decide dónde ubicarlo y lo mantiene.
\end{ejemplobox}
\begin{borradorbox}{máx. 560 caracteres (≈ 4 renglones)}
\renglones{4}
\end{borradorbox}

% =====================================================================
\section{Área de impacto y ODS: elige con coherencia}

El formulario pide elegir \textbf{hasta 2 áreas de impacto}, y según ellas se habilitan los \textbf{Objetivos de Desarrollo Sostenible (ODS)} relacionados (entre 1 y 4). Elige los que tu proyecto \textbf{de verdad} aborda: la coherencia suma; el relleno resta.

\begin{center}
\renewcommand{\arraystretch}{1.3}
\begin{tabularx}{\linewidth}{@{}p{0.30\linewidth} X@{}}
\toprule
\textbf{\color{smjAzul}Área de impacto} & \textbf{\color{smjAzul}ODS que se habilitan}\\
\midrule
\textbf{Sostenibilidad ambiental} & 6 Agua limpia · 7 Energía · 12 Producción responsable · 13 Clima · 14 Océanos · 15 Vida terrestre\\
\textbf{Salud y bienestar} & 2 Hambre cero · 3 Salud y bienestar · 6 Agua limpia\\
\textbf{Apoyo a la comunidad local} & 1 Fin de la pobreza · 4 Educación · 5 Igualdad de género · 10 Desigualdades · 11 Ciudades sostenibles · 16 Paz y justicia · 17 Alianzas\\
\textbf{Deportes} & 3 Salud y bienestar · 4 Educación · 10 Desigualdades · 11 Ciudades sostenibles\\
\bottomrule
\end{tabularx}
\end{center}

\noindent\textbf{En el ejemplo:} AquaVigía marcaría \emph{Sostenibilidad ambiental} y \emph{Salud y bienestar}, con los ODS \textbf{6} (agua limpia y saneamiento) y \textbf{3} (salud y bienestar). Encajan sin forzar nada.

% =====================================================================
\section{STEM con sentido: que se vea funcionando}

Este es el sello del concurso. En la etapa 4 debes mostrar el STEM \emph{actuando} dentro de tu solución. Usa esta tabla como chequeo: trata de poder llenar al menos dos columnas con honestidad.

\begin{center}
\renewcommand{\arraystretch}{1.3}
\begin{tabularx}{\linewidth}{@{}p{0.20\linewidth} X@{}}
\toprule
\textbf{\color{smjAzul}Letra} & \textbf{\color{smjAzul}Pregúntate\ldots}\\
\midrule
\textbf{C — Ciencia} & ¿Qué principio o conocimiento natural explica por qué funciona? (un rango seguro, una reacción, un fenómeno físico)\\
\textbf{T — Tecnología} & ¿Qué herramienta digital o dispositivo usas? (sensor, app, microcontrolador, hoja de cálculo, IA)\\
\textbf{I — Ingeniería} & ¿Qué diseñas o construyes y cómo lo haces resistente, usable y reparable?\\
\textbf{M — Matemática} & ¿Qué cálculo, dato o umbral sostiene tus decisiones? (promedios, porcentajes, comparaciones)\\
\bottomrule
\end{tabularx}
\end{center}

% =====================================================================
\section{Si llegan a la final: el pitch}

Los equipos destacados presentan un \textbf{pitch} (exposición breve) ante el jurado. Un buen pitch cuenta una historia en este orden, en pocos minutos:

\begin{enumerate}[label=\textbf{\arabic*.}]
  \item \textbf{El problema, con rostro.} Empieza por la comunidad y una escena concreta. Que el jurado \emph{sienta} el problema.
  \item \textbf{La solución, en una frase.} Qué es y qué hace, sin tecnicismos.
  \item \textbf{El STEM, mostrando el cómo.} Aquí sí, el conocimiento funcionando.
  \item \textbf{El impacto, con un número.} Cómo sabrán que cambió algo.
  \item \textbf{El equipo y el siguiente paso.} Quiénes son y qué harían con el apoyo.
\end{enumerate}
\noindent\textbf{Consejo:} ensayen con el reloj, repartan las voces del equipo y preparen respuestas a «¿por qué su idea y no otra?» y «¿cómo lo sostendrían en el tiempo?».

% =====================================================================
\section{Rúbrica de autoevaluación «alto valor»}

Antes de enviar, califiquen cada criterio con sinceridad. Si algo cae en \emph{En desarrollo}, aún tienen tiempo de subirlo.

\begin{center}
\renewcommand{\arraystretch}{1.35}
\begin{tabularx}{\linewidth}{@{}p{0.22\linewidth} X X@{}}
\toprule
\textbf{\color{smjAzul}Criterio} & \textbf{\color{smjVino}En desarrollo} & \textbf{\color{valorVerde}Alto valor}\\
\midrule
Empatía & Describimos a la comunidad de oídas. & Hablamos con ella; tenemos voces y datos.\\
Problema & Es amplio o es «la falta de nuestra idea». & Es concreto, ubicado y con impacto claro.\\
Solución & Genérica; no decimos en qué nos diferenciamos. & Clara en una frase y diferenciada.\\
STEM & Se menciona como etiqueta. & Se ve funcionando en el prototipo.\\
Impacto & «Será de gran ayuda». & Decimos cómo lo mediremos, con la comunidad.\\
\bottomrule
\end{tabularx}
\end{center}

% =====================================================================
\section{La lente MILC: dale profundidad a tu propuesta}

Un proyecto Solve for Tomorrow crece cuando lo miras con las \textbf{cinco dimensiones} del modelo MILC. Úsalas para encontrar y enriquecer tu idea:

\begin{milcbox}
\begin{itemize}[nosep]
  \item \textbf{Personal:} ¿qué de este problema te toca a ti o a tu familia? Ahí hay energía verdadera.
  \item \textbf{Emocional:} ¿qué sienten quienes lo viven? Esa emoción es el corazón de tu pitch.
  \item \textbf{Ciudadana:} ¿qué derecho está en juego? (agua, salud, educación). Eso le da peso.
  \item \textbf{Local:} ¿por qué tu solución encaja en Cartago y no es copia de otro lugar? El territorio es tu ventaja.
  \item \textbf{Intergeneracional:} ¿quién más se beneficia, ahora y después? Las soluciones que duran convencen.
\end{itemize}
\end{milcbox}

% =====================================================================
\section{Checklist de inscripción (en el sitio oficial)}

Cuando la propuesta esté pulida, el \textbf{docente} la inscribe en el sitio oficial. Ten a mano:

\begin{multicols}{2}
\textbf{\color{smjAzul}1. Datos del docente}\\
\chk Correo (máx. 100) \quad \chk ID de profesor (máx. 10)\\
\chk Nombre completo \quad \chk Sexo \quad \chk Celular

\smallskip
\textbf{\color{smjAzul}2. Datos de la institución}\\
\chk Nombre del colegio \quad \chk Dirección\\
\chk País \quad \chk Departamento \quad \chk Ciudad

\smallskip
\textbf{\color{smjAzul}3. Datos del equipo}\\
\chk Nombre del equipo \quad \chk Cuántos integrantes (2–4)\\
\chk Por estudiante: grado y sexo

\smallskip
\textbf{\color{smjAzul}4. La propuesta}\\
\chk Área(s) de impacto (máx. 2) \quad \chk ODS (1–4)\\
\chk Empatizar \quad \chk Definir \quad \chk Idear\\
\chk Prototipar (STEM) \quad \chk Validar

\smallskip
\textbf{\color{smjAzul}5. Cierre}\\
\chk ¿Participaron antes? \quad \chk Tratamiento de datos\\
\chk Términos y condiciones (los lee y acepta el docente)
\end{multicols}

\begin{oficial}
\textbf{La inscripción es exclusivamente oficial.} Esta guía no inscribe proyectos. Verifica fechas, bases y formulario en:\\
\url{https://www.samsung.com/co/campaign/solvefortomorrow/}
\end{oficial}

\vfill
\begin{center}
\footnotesize\color{gray}
\rule{0.4\linewidth}{0.4pt}\\[4pt]
\textit{Solve for Tomorrow} es un programa y marca de Samsung Electronics. Esta es una guía de preparación pedagógica \textbf{independiente}, elaborada por la I.E. Sor María Juliana (Cartago) en su plataforma ConectaTE, \textbf{sin afiliación ni respaldo oficial} de Samsung.\\[3pt]
PhD. Álvaro Cárdenas Orozco \;·\; Área de Tecnología e Informática \;·\; Cartago, Valle del Cauca \;·\; 2026
\end{center}

\end{document}
